\documentclass[10pt,a4paper]{amsart}
\usepackage{amsmath,amssymb,amsthm,graphicx}
% Retained preamble for the Stacks Project chapter snapshot in this directory.
% The chapter's own preamble.tex is a copy of the shared Stacks Project preamble; the block
% below is sliced MECHANICALLY, line for line, from it. Nothing was added, removed or
% redefined; the wrapper class, the xr-hyper/hyperref wiring and the externaldocument list
% are omitted because the assembler supplies the document class.
%
% lines taken from preamble.tex: 3, 5, 9, 11, 12, 156, 157, 158, ... (51 lines total)
%
% package coverage restored from the chapter's own preamble (2026-10-13): xy, xy-2cell

\usepackage[all]{xy}
\xyoption{2cell}
\UseAllTwocells

\newenvironment{reference}{\comment}{\endcomment}
\newenvironment{slogan}{\comment}{\endcomment}
\newenvironment{history}{\comment}{\endcomment}
\theoremstyle{plain}
\newtheorem{theorem}[subsection]{Theorem}
\newtheorem{proposition}[subsection]{Proposition}
\newtheorem{lemma}[subsection]{Lemma}
\theoremstyle{definition}
\newtheorem{definition}[subsection]{Definition}
\newtheorem{example}[subsection]{Example}
\newtheorem{exercise}[subsection]{Exercise}
\newtheorem{situation}[subsection]{Situation}
\theoremstyle{remark}
\newtheorem{remark}[subsection]{Remark}
\newtheorem{remarks}[subsection]{Remarks}
\numberwithin{equation}{subsection}
\def\lim{\mathop{\mathrm{lim}}\nolimits}
\def\colim{\mathop{\mathrm{colim}}\nolimits}
\def\Spec{\mathop{\mathrm{Spec}}}
\def\Hom{\mathop{\mathrm{Hom}}\nolimits}
\def\Ext{\mathop{\mathrm{Ext}}\nolimits}
\def\SheafHom{\mathop{\mathcal{H}\!\mathit{om}}\nolimits}
\def\SheafExt{\mathop{\mathcal{E}\!\mathit{xt}}\nolimits}
\def\Sch{\mathit{Sch}}
\def\Mor{\mathop{\mathrm{Mor}}\nolimits}
\def\Ob{\mathop{\mathrm{Ob}}\nolimits}
\def\Sh{\mathop{\mathit{Sh}}\nolimits}
\def\NL{\mathop{N\!L}\nolimits}
\def\CH{\mathop{\mathrm{CH}}\nolimits}
\def\proetale{{pro\text{-}\acute{e}tale}}
\def\etale{{\acute{e}tale}}
\def\QCoh{\mathit{QCoh}}
\def\Ker{\mathop{\mathrm{Ker}}}
\def\Im{\mathop{\mathrm{Im}}}
\def\Coker{\mathop{\mathrm{Coker}}}
\def\Coim{\mathop{\mathrm{Coim}}}
\DeclareMathSymbol{\boxtimes}{\mathbin}{AMSa}{"02}
\def\QCohstack{\mathcal{QC}\!\mathit{oh}}
\def\Cohstack{\mathcal{C}\!\mathit{oh}}
\def\Spacesstack{\mathcal{S}\!\mathit{paces}}
\def\Quotfunctor{\mathrm{Quot}}
\def\Hilbfunctor{\mathrm{Hilb}}
\def\Curvesstack{\mathcal{C}\!\mathit{urves}}
\def\Polarizedstack{\mathcal{P}\!\mathit{olarized}}
\def\Complexesstack{\mathcal{C}\!\mathit{omplexes}}
\def\Pic{\mathop{\mathrm{Pic}}\nolimits}
\def\Picardstack{\mathcal{P}\!\mathit{ic}}
\def\Picardfunctor{\mathrm{Pic}}
\def\Deformationcategory{\mathcal{D}\!\mathit{ef}}
\usepackage{cleveref}
\crefname{originallemma}{original lemma}{original lemmas}
\Crefname{originallemma}{Original lemma}{Original lemmas}
\makeatletter
\newlabel{homology-lemma-double-complex-gives-resolution}{{0133}{source}{}{sourcepointer-1}{}}
\newlabel{homology-lemma-double-complex-gives-resolution@cref}{{[originallemma][0][]0133}{[1][0][]source}{}{}{}}
\makeatother
\crefname{originallemma}{original lemma}{original lemmas}
\Crefname{originallemma}{Original lemma}{Original lemmas}
\makeatletter
\newlabel{homology-section-double-complex}{{012X}{source}{}{sourcepointer-2}{}}
\newlabel{homology-section-double-complex@cref}{{[originallemma][0][]012X}{[1][0][]source}{}{}{}}
\makeatother
\crefname{originallemma}{original lemma}{original lemmas}
\Crefname{originallemma}{Original lemma}{Original lemmas}
\makeatletter
\newlabel{lemma-fibre-products-preserve-properties}{{lemma-fibre-products-preserve-properties}{source}{}{sourcepointer-3}{}}
\newlabel{lemma-fibre-products-preserve-properties@cref}{{[originallemma][0][]lemma-fibre-products-preserve-properties}{[1][0][]source}{}{}{}}
\makeatother
\crefname{originallemma}{original lemma}{original lemmas}
\Crefname{originallemma}{Original lemma}{Original lemmas}
\makeatletter
\newlabel{lemma-pullback-injective}{{lemma-pullback-injective}{source}{}{sourcepointer-4}{}}
\newlabel{lemma-pullback-injective@cref}{{[originallemma][0][]lemma-pullback-injective}{[1][0][]source}{}{}{}}
\makeatother
\crefname{originallemma}{original lemma}{original lemmas}
\Crefname{originallemma}{Original lemma}{Original lemmas}
\makeatletter
\newlabel{lemma-pushforward-injective}{{lemma-pushforward-injective}{source}{}{sourcepointer-5}{}}
\newlabel{lemma-pushforward-injective@cref}{{[originallemma][0][]lemma-pushforward-injective}{[1][0][]source}{}{}{}}
\makeatother
\crefname{originallemma}{original lemma}{original lemmas}
\Crefname{originallemma}{Original lemma}{Original lemmas}
\makeatletter
\newlabel{proposition-exactness-cech-complex}{{proposition-exactness-cech-complex}{source}{}{sourcepointer-6}{}}
\newlabel{proposition-exactness-cech-complex@cref}{{[originallemma][0][]proposition-exactness-cech-complex}{[1][0][]source}{}{}{}}
\makeatother
\title{Stacks, Sheaves on Algebraic Stacks: lemma-cech-to-cohomology}
\author{The Stacks Project Authors}
\date{stacks-project:stacks-sheaves.tex:local-snapshot-20261012; GNU FDL 1.2 or later}
\begin{document}
\maketitle
\noindent\textit{Editorial note.} Statement span is lines 3600-3621 of sources/stacks-algebraic-stacks-sheaves/pinned-original-source.tex; proof span is lines 3623-3669, the complete author proof (45 body lines).\par\medskip
\noindent\textit{Editorial note.} Navigation stubs: 6 cited-but-not-included labels are supplied as non-load-bearing pointers with their pinned tags where the snapshot tag index has them.\par\medskip
\noindent\textit{Editorial note.} extra\_packages is empty on purpose: the retained preamble already loads xy with options and a second declaration clashes in the assembler wrapper.\par\medskip
\noindent\textit{Editorial note.} Text is retained verbatim, including the author's own typos and placeholders.\par\medskip
\noindent\textit{Editorial note.} Queue index 12 of queue\_algebraic-stacks-sheaves.json (pb=1973, stmt 3600-3621, proof 3623-3669).\par\medskip
\noindent\textit{Allowed standard prerequisites (editorial).}\begin{itemize}
\item Explicit prerequisite boundary. The cited inputs are prerequisites; the argument in the retained span is not replaced by any cited result.
\end{itemize}\medskip
\ifdefined\hypertarget\hypertarget{sourcepointer-1}{}\fi
\noindent\textit{Source reference.} Original lemma 0133: sources/stacks-algebraic-stacks-sheaves/pinned-original-source.tex, source lines 3669--3669. Cited by the retained statement or proof but not part of the collected spans; supplied only as a navigation pointer. The argument does not depend on its content.\par\medskip
\ifdefined\hypertarget\hypertarget{sourcepointer-2}{}\fi
\noindent\textit{Source reference.} Original lemma 012X: sources/stacks-algebraic-stacks-sheaves/pinned-original-source.tex, source lines 3669--3669. Cited by the retained statement or proof but not part of the collected spans; supplied only as a navigation pointer. The argument does not depend on its content.\par\medskip
\ifdefined\hypertarget\hypertarget{sourcepointer-3}{}\fi
\noindent\textit{Source reference.} Original lemma lemma-fibre-products-preserve-properties: sources/stacks-algebraic-stacks-sheaves/pinned-original-source.tex, source lines 2945--2954. Cited by the retained statement or proof but not part of the collected spans; supplied only as a navigation pointer. The argument does not depend on its content.\par\medskip
\ifdefined\hypertarget\hypertarget{sourcepointer-4}{}\fi
\noindent\textit{Source reference.} Original lemma lemma-pullback-injective: sources/stacks-algebraic-stacks-sheaves/pinned-original-source.tex, source lines 3009--3022. Cited by the retained statement or proof but not part of the collected spans; supplied only as a navigation pointer. The argument does not depend on its content.\par\medskip
\ifdefined\hypertarget\hypertarget{sourcepointer-5}{}\fi
\noindent\textit{Source reference.} Original lemma lemma-pushforward-injective: sources/stacks-algebraic-stacks-sheaves/pinned-original-source.tex, source lines 2835--2846. Cited by the retained statement or proof but not part of the collected spans; supplied only as a navigation pointer. The argument does not depend on its content.\par\medskip
\ifdefined\hypertarget\hypertarget{sourcepointer-6}{}\fi
\noindent\textit{Source reference.} Original lemma proposition-exactness-cech-complex: sources/stacks-algebraic-stacks-sheaves/pinned-original-source.tex, source lines 3558--3568. Cited by the retained statement or proof but not part of the collected spans; supplied only as a navigation pointer. The argument does not depend on its content.\par\medskip
\begin{lemma}
\label{lemma-cech-to-cohomology}
Let $f : \mathcal{U} \to \mathcal{X}$ be a $1$-morphism of categories fibred
in groupoids over $(\Sch/S)_{fppf}$. Let
$\tau \in \{Zar, \etale, smooth, syntomic, fppf\}$.
Assume
\begin{enumerate}
\item $\mathcal{F}$ is an abelian sheaf on $\mathcal{X}_\tau$,
\item for every object $x$ of $\mathcal{X}$ there exists a covering
$\{x_i \to x\}$ in $\mathcal{X}_\tau$ such that each $x_i$ is isomorphic
to $f(u_i)$ for some object $u_i$ of $\mathcal{U}$,
\item the category $\mathcal{U}$ has equalizers, and
\item the functor $f$ is faithful.
\end{enumerate}
Then there is a first quadrant spectral sequence of abelian groups
$$
E_1^{p, q} = H^q((\mathcal{U}_p)_\tau, f_p^{-1}\mathcal{F})
\Rightarrow
H^{p + q}(\mathcal{X}_\tau, \mathcal{F})
$$
converging to the cohomology of $\mathcal{F}$ in the $\tau$-topology.
\end{lemma}

\begin{proof}
Before we start the proof we make some remarks. By
Lemma \ref{lemma-fibre-products-preserve-properties}
(and induction) all of the categories fibred in groupoids $\mathcal{U}_p$
have equalizers and all of the morphisms $f_p : \mathcal{U}_p \to \mathcal{X}$
are faithful. Let $\mathcal{I}$ be an injective object
of $\textit{Ab}(\mathcal{X}_\tau)$. By
Lemma \ref{lemma-pullback-injective}
we see $f_p^{-1}\mathcal{I}$ is an injective object of
$\textit{Ab}((\mathcal{U}_p)_\tau)$.
Hence $f_{p, *}f_p^{-1}\mathcal{I}$ is an injective object of
$\textit{Ab}(\mathcal{X}_\tau)$ by
Lemma \ref{lemma-pushforward-injective}.
Hence
Proposition \ref{proposition-exactness-cech-complex}
shows that the extended relative {\v C}ech complex
$$
\ldots \to 0 \to
\mathcal{I} \to
f_{0, *}f_0^{-1}\mathcal{I} \to
f_{1, *}f_1^{-1}\mathcal{I} \to
f_{2, *}f_2^{-1}\mathcal{I} \to \ldots
$$
is an exact complex in $\textit{Ab}(\mathcal{X}_\tau)$ all of whose
terms are injective. Taking global sections of this complex is exact
and we see that the {\v C}ech complex
$\check{\mathcal{C}}^\bullet(\mathcal{U} \to \mathcal{X}, \mathcal{I})$
is quasi-isomorphic to $\Gamma(\mathcal{X}_\tau, \mathcal{I})[0]$.

\medskip\noindent
With these preliminaries out of the way consider the two spectral sequences
associated to the double complex (see
Homology, Section \ref{homology-section-double-complex})
$$
\check{\mathcal{C}}^\bullet(\mathcal{U} \to \mathcal{X}, \mathcal{I}^\bullet)
$$
where $\mathcal{F} \to \mathcal{I}^\bullet$ is an injective resolution
in $\textit{Ab}(\mathcal{X}_\tau)$.
The discussion above shows that
Homology, Lemma \ref{homology-lemma-double-complex-gives-resolution}
applies which shows that
$\Gamma(\mathcal{X}_\tau, \mathcal{I}^\bullet)$
is quasi-isomorphic to the total complex associated to the double complex.
By our remarks above the complex $f_p^{-1}\mathcal{I}^\bullet$ is an
injective resolution of $f_p^{-1}\mathcal{F}$. Hence the other spectral
sequence is as indicated in the lemma.
\end{proof}

\end{document}
